% GENERATED FILE. Edit canonical lessons, then run npm run generate:books.
% canonical-source-hash: fnv1a64:9fe8857f4dbdfb96
% canonical-lessons: IT-C514-tuffo, IT-C514-cuffia, IT-C514-vela, IT-C514-immergersi, IT-C514-trampolino

\chapter{In acqua}
\label{ch:it-in-acqua}
\hlchaptermodality{514}

\begin{chapteropening}
\textbf{By the end of this chapter:} \emph{I can talk about water sports.}

Complete the last lesson of chapter 514: I can talk about water sports.
\end{chapteropening}

\section[il tuffo]{il tuffo — dive, plunge}

\label{lesson:IT-C514-tuffo}



\begin{quote}
Before the new one: say the Italian for athlete, then the Italian for to pant.
\end{quote}

\subsection*{You'll want to know: il tuffo}
\textbf{il tuffo} — \textquotedblleft{}dive, plunge\textquotedblright{}.

\begin{quote}
Ho fatto un tuffo dal molo.
\end{quote}

\begin{quote}
\emph{I dived in from the pier.}
\end{quote}

\emph{Fare un tuffo} — to go for a dip.

\begin{cousinweb}
The origin of \emph{tuffo} is not settled; it is often taken to imitate the sound of a body hitting water, or traced to a Germanic word for dipping.
\end{cousinweb}

\begin{grammarlens}[title={What you've built}]
One more word for telling, explaining and describing at B1.
\end{grammarlens}

\subsection*{Your turn}
\begin{itemize}
\raggedright
  \item \emph{Say it:} \emph{il tuffo}
  \item \emph{Say it:} \emph{il tuffo}, once more
  \item \emph{Say it:} a sentence of your own with \emph{il tuffo}
  \item \emph{Recall:} say the Italian for athlete, then the Italian for to pant, then say \emph{il tuffo} again
\end{itemize}

\begin{tcolorbox}[breakable,colback=teal!4,colframe=teal!35!black,title={Before you move on}]
What does il tuffo mean? (\textquotedblleft{}Dive, plunge\textquotedblright{}.) Say it once more.
\end{tcolorbox}
\section[la cuffia]{la cuffia — swimming cap; headphones}

\label{lesson:IT-C514-cuffia}



\begin{quote}
Before the new one: say the Italian for to pant, then the Italian for dive.
\end{quote}

\subsection*{You'll want to know: la cuffia}
\textbf{la cuffia} — \textquotedblleft{}swimming cap; headphones\textquotedblright{}.

\begin{quote}
In piscina la cuffia è obbligatoria.
\end{quote}

\begin{quote}
\emph{In the pool a swimming cap is compulsory.}
\end{quote}

Also headphones: \emph{le cuffie}.

\begin{cousinweb}
\emph{cuffia} comes from Late Latin \emph{cofia} \textquotedblleft{}a cap, a bonnet\textquotedblright{} (EtymDB), probably of Germanic origin.
\end{cousinweb}

\begin{grammarlens}[title={What you've built}]
One more word for telling, explaining and describing at B1.
\end{grammarlens}

\subsection*{Your turn}
\begin{itemize}
\raggedright
  \item \emph{Say it:} \emph{la cuffia}
  \item \emph{Say it:} \emph{la cuffia}, once more
  \item \emph{Say it:} a sentence of your own with \emph{la cuffia}
  \item \emph{Recall:} say the Italian for to pant, then the Italian for dive, then say \emph{la cuffia} again
\end{itemize}

\begin{tcolorbox}[breakable,colback=teal!4,colframe=teal!35!black,title={Before you move on}]
What does la cuffia mean? (\textquotedblleft{}Swimming cap; headphones\textquotedblright{}.) Say it once more.
\end{tcolorbox}
\section[la vela]{la vela — sail; sailing}

\label{lesson:IT-C514-vela}



\begin{quote}
Before the new one: say the Italian for dive, then the Italian for swimming cap.
\end{quote}

\subsection*{You'll want to know: la vela}
\textbf{la vela} — \textquotedblleft{}sail; sailing\textquotedblright{}.

\begin{quote}
D'estate faccio un corso di vela.
\end{quote}

\begin{quote}
\emph{In summer I take a sailing course.}
\end{quote}

\emph{Andare a vela} — to sail; \emph{la barca a vela} — sailing boat.

\begin{cousinweb}
\emph{vela} is inherited from Latin \emph{vēla}, the plural of \emph{vēlum} \textquotedblleft{}a sail, a covering\textquotedblright{} (EtymDB).
\end{cousinweb}

\begin{grammarlens}[title={What you've built}]
One more word for telling, explaining and describing at B1.
\end{grammarlens}

\subsection*{Your turn}
\begin{itemize}
\raggedright
  \item \emph{Say it:} \emph{la vela}
  \item \emph{Say it:} \emph{la vela}, once more
  \item \emph{Say it:} a sentence of your own with \emph{la vela}
  \item \emph{Recall:} say the Italian for dive, then the Italian for swimming cap, then say \emph{la vela} again
\end{itemize}

\begin{tcolorbox}[breakable,colback=teal!4,colframe=teal!35!black,title={Before you move on}]
What does la vela mean? (\textquotedblleft{}Sail; sailing\textquotedblright{}.) Say it once more.
\end{tcolorbox}
\section[immergersi]{immergersi — to dive under, to immerse oneself}

\label{lesson:IT-C514-immergersi}



\begin{quote}
Before the new one: say the Italian for swimming cap, then the Italian for sail.
\end{quote}

\subsection*{You'll want to know: immergersi}
\textbf{immergersi} — \textquotedblleft{}to dive under, to immerse oneself\textquotedblright{}.

\begin{quote}
Ci siamo immersi per vedere i pesci.
\end{quote}

\begin{quote}
\emph{We dived under to watch the fish.}
\end{quote}

The past participle is \emph{immerso}. \emph{Immergersi nella lettura} — to lose oneself in reading.

\begin{cousinweb}
\emph{immergersi} comes from Latin \emph{immergere} \textquotedblleft{}to dip into\textquotedblright{} (Lewis \& Short), from \emph{mergere} \textquotedblleft{}to dip, to sink\textquotedblright{}.
\end{cousinweb}

\begin{grammarlens}[title={What you've built}]
One more word for telling, explaining and describing at B1.
\end{grammarlens}

\subsection*{Your turn}
\begin{itemize}
\raggedright
  \item \emph{Say it:} \emph{immergersi}
  \item \emph{Say it:} \emph{immergersi}, once more
  \item \emph{Say it:} a sentence of your own with \emph{immergersi}
  \item \emph{Recall:} say the Italian for swimming cap, then the Italian for sail, then say \emph{immergersi} again
\end{itemize}

\begin{tcolorbox}[breakable,colback=teal!4,colframe=teal!35!black,title={Before you move on}]
What does immergersi mean? (\textquotedblleft{}To dive under, to immerse oneself\textquotedblright{}.) Say it once more.
\end{tcolorbox}
\section[il trampolino]{il trampolino — diving board; springboard}

\label{lesson:IT-C514-trampolino}



\begin{quote}
Before the new one: say the Italian for sail, then the Italian for to dive under.
\end{quote}

\subsection*{You'll want to know: il trampolino}
\textbf{il trampolino} — \textquotedblleft{}diving board; springboard\textquotedblright{}.

\begin{quote}
Ha paura di saltare dal trampolino alto.
\end{quote}

\begin{quote}
\emph{He is afraid to jump from the high board.}
\end{quote}

Also a ski jump.

\begin{cousinweb}
\emph{trampolino} is built on \emph{trampolo} \textquotedblleft{}stilt\textquotedblright{} (EtymDB), which comes from a German word like \emph{trampeln} \textquotedblleft{}to trample\textquotedblright{}.
\end{cousinweb}

\begin{grammarlens}[title={What you've built}]
One more word for telling, explaining and describing at B1.
\end{grammarlens}

\subsection*{Your turn}
\begin{itemize}
\raggedright
  \item \emph{Say it:} \emph{il trampolino}
  \item \emph{Say it:} \emph{il trampolino}, once more
  \item \emph{Say it:} a sentence of your own with \emph{il trampolino}
  \item \emph{Recall:} say the Italian for sail, then the Italian for to dive under, then say \emph{il trampolino} again
\end{itemize}

\begin{tcolorbox}[breakable,colback=teal!4,colframe=teal!35!black,title={Before you move on}]
What does il trampolino mean? (\textquotedblleft{}Diving board; springboard\textquotedblright{}.) Say it once more.
\end{tcolorbox}
